Estas notas foram elaboradas com o objetivo de acompanhar o curso de Análise no $\mathbb{R}^n$, reunindo de forma organizada os principais conceitos, resultados e técnicas da análise em várias variáveis. O propósito é servir como material de estudo e referência, apresentando os tópicos do curso com rigor matemático e, ao mesmo tempo, com uma exposição clara e progressiva.

O curso inicia com o estudo da diferenciabilidade de aplicações entre espaços euclidianos, incluindo resultados fundamentais como a regra da cadeia e a desigualdade do valor médio. Em seguida, são desenvolvidos teoremas centrais da teoria local de aplicações diferenciáveis, como o teorema da função inversa e o teorema da função implícita, bem como a forma local de submersões e imersões e o teorema do posto.

Posteriormente, são introduzidos conceitos relacionados à teoria da medida e à integração em várias variáveis, incluindo conjuntos de medida nula e mudanças de variáveis em integrais múltiplas. Esses resultados constituem ferramentas essenciais para o estudo geométrico e analítico de aplicações diferenciáveis.

Na parte final do curso, são estudadas formas alternadas e formas diferenciais, juntamente com a diferencial exterior e partições da unidade. Esses conceitos permitem formular e demonstrar resultados fundamentais da análise e da geometria diferencial, como integrais de superfície e o Teorema de Stokes, culminando no estudo de propriedades de invariância por homotopia.

Estas notas foram preparadas com base em referências clássicas da área, em particular os textos de Elon Lages Lima, Michael Spivak e outros autores que contribuíram significativamente para o desenvolvimento da análise em várias variáveis.

\section*{Ementa}

\begin{itemize}
    \item Diferenciabilidade de uma aplicação; regra da cadeia; desigualdade do valor médio; teorema da aplicação implícita e inversa; forma local das submersões e das imersões; teorema do posto;
    
    \item conjuntos de medida nula; mudanças de variáveis em integrais múltiplas;
    
    \item formas alternadas; formas diferenciais; a diferencial exterior; partição da unidade;
    
    \item integrais de superfície; teorema de Stokes; integral invariante por homotopia.
\end{itemize}

\section*{Bibliografia}

\begin{enumerate}
    \item Lima, E.L. \textit{Curso de Análise}, vol. 2. Projeto Euclides, SBM, 1981.
    
    \item Lima, E.L. \textit{Análise no Espaço $\mathbb{R}^n$}. Edgar Blücher, 1970.
    
    \item Spivak, M. \textit{Calculus on Manifolds}. W. A. Benjamin, New York, 1965.
    
    \item Giaquinta, M.; Modica, G. \textit{Mathematical Analysis: An Introduction to Functions of Several Variables}.
\end{enumerate}
